% auto-generated -- do not edit by hand
\begin{longtable}{@{}p{0.9cm} p{2.9cm} p{1.5cm} p{8.4cm}@{}}
\caption{The 60-paper corpus. Folder A = \texttt{Agents/}, C = \texttt{Constraints/}.}
\label{tab:corpus}\\
\toprule
Folder & First author & arXiv & Title \\
\midrule
\endfirsthead
\toprule
Folder & First author & arXiv & Title \\
\midrule
\endhead
\bottomrule
\endfoot
A & Furlanetto et al. & \texttt{astro-ph/0312435} & Large-Scale Structure Shocks at Low and High Redshifts \\
A & Furlanetto et al. & \texttt{0910.4410} & Secondary ionization and heating by fast electrons \\
A & Kulkarni et al. & \texttt{1310.0684} & Chemical constraints on the contribution of Population III stars to cosmic reionization \\
A & Pacucci et al. & \texttt{1403.6125} & The X-ray spectra of the first galaxies: 21cm signatures \\
A & Tueros et al. & \texttt{1409.6225} & Cosmic reionization by primordial cosmic rays \\
A & Robertson et al. & \texttt{1502.02024} & Cosmic Reionization and Early Star-Forming Galaxies: A Joint Analysis of New Constraints from Planck and Hubble Space Telescope \\
A & Madau et al. & \texttt{1507.07678} & Cosmic Reionization after Planck: Could Quasars Do It All? \\
A & Sazonov et al. & \texttt{1509.08408} & Preheating of the Universe by cosmic rays from primordial supernovae at the beginning of cosmic reionization \\
A & Madau et al. & \texttt{1606.07887} & Radiation Backgrounds at Cosmic Dawn: X-Rays from Compact Binaries \\
A & Rosdahl et al. & \texttt{1801.07259} & The SPHINX Cosmological Simulations of the First Billion Years: the Impact of Binary Stars on Reionization \\
A & Sharma & \texttt{1804.05843} & Astrophysical radio background cannot explain the EDGES 21-cm signal: constraints from cooling of non-thermal electrons \\
A & Kulkarni et al. & \texttt{1807.09774} & Evolution of the AGN UV luminosity function from redshift 7.5 \\
A & Finkelstein et al. & \texttt{1902.02792} & Conditions for Reionizing the Universe with A Low Galaxy Ionizing Photon Escape Fraction \\
A & Fialkov et al. & \texttt{1902.02438} & Signature of Excess Radio Background in the 21-cm Global Signal and Power Spectrum \\
A & Naidu et al. & \texttt{1907.13130} & Rapid Reionization by the Oligarchs: The Case for Massive, UV-Bright, Star-Forming Galaxies with High Escape Fractions \\
A & Dayal et al. & \texttt{2001.06021} & Reionization with galaxies and active galactic nuclei \\
A & D'Aloisio et al. & \texttt{2002.02467} & Hydrodynamic Response of the Intergalactic Medium to Reionization \\
A & Reis et al. & \texttt{2008.04315} & High-redshift radio galaxies: a potential new source of 21-cm fluctuations \\
A & Eide et al. & \texttt{2009.06631} & Large scale simulations of H and He reionization and heating driven by stars and more energetic sources \\
A & Reis et al. & \texttt{2101.01777} & The subtlety of Ly-a photons: changing the expected range of the 21-cm signal \\
A & Xu et al. & \texttt{2102.12865} & Maximum Absorption of the Global 21 cm Spectrum in the Standard Cosmological Model \\
A & Murphy et al. & \texttt{2105.06900} & Ionizing photon production of Population III stars: effects of rotation, convection, and initial mass function \\
A & Kannan et al. & \texttt{2110.00584} & Introducing the THESAN project: radiation-magnetohydrodynamic simulations of the epoch of reionization \\
A & Wang et al. & \texttt{2212.04476} & A strong He II $\lambda$1640 emitter with extremely blue UV spectral slope at z=8.16: presence of Pop III stars? \\
A & Gessey-Jones et al. & \texttt{2304.07201} & Signatures of Cosmic Ray Heating in 21-cm Observables \\
A & Atek et al. & \texttt{2308.08540} & Most of the photons that reionized the Universe came from dwarf galaxies \\
A & Munoz et al. & \texttt{2404.07250} & Reionization after JWST: a photon budget crisis? \\
A & Simmonds et al. & \texttt{2409.01286} & Ionising properties of galaxies in JADES for a stellar mass complete sample: resolving the cosmic ionising photon budget crisis at \\
A & Asthana et al. & \texttt{2409.15453} & The impact of faint AGN discovered by JWST on reionization \\
A & Wasserman et al. & \texttt{2507.21764} & Ultraviolet photon production rates of the first stars: Impact on the He II $\lambda$ 1640 A emission line from primordial star cl \\
A & Singha et al. & \texttt{2510.26990} & Faint active galactic nuclei supplied 31-75\% of hydrogen-ionizing photons at z>5 \\
A & Austin et al. & \texttt{2512.10839} & Resolving the ionizing photon budget crisis with JWST/NIRCam HII clumping constraints at z=6 \\
A & Maiolino et al. & \texttt{2603.20362} & The search for Population III: Confirmation of a HeII emitter with no metal lines at z=10.6 \\
A & Dasgupta et al. & \texttt{2604.11542} & Impact of Stochastic Pop III X-ray Binaries on the Cosmological 21-cm Signal \\
A & Giovinazzo et al. & \texttt{2607.22834} & Reionization driven by the few: the ionizing budget of galaxies at z=5-10 from JWST/NIRSpec \\
C & Becker et al. & \texttt{1407.4850} & Evidence of patchy hydrogen reionization from an extreme Ly$\alpha$ trough below redshift six \\
C & McGreer et al. & \texttt{1411.5375} & Model-independent evidence in favor of an end to reionization by z 6 \\
C & Mason et al. & \texttt{1709.05356} & The Universe is Reionizing at z 7: Bayesian Inference of the IGM Neutral Fraction Using Ly$\alpha$ Emission from Galaxies \\
C & Davies et al. & \texttt{1802.06066} & Quantitative Constraints on the Reionization History from the IGM Damping Wing Signature in Two Quasars at z > 7 \\
C & HERA Collab. et al. & \texttt{1807.06209} & Planck 2018 results. VI. Cosmological parameters \\
C & Pagano et al. & \texttt{1908.09856} & Reionization optical depth determination from Planck HFI data with ten percent accuracy \\
C & Gaikwad et al. & \texttt{2001.10018} & Probing the thermal state of the intergalactic medium at z>5 with the transmission spikes in high-resolution Ly$\alpha$ forest spe \\
C & Reichardt et al. & \texttt{2002.06197} & An Improved Measurement of the Secondary Cosmic Microwave Background Anisotropies from the SPT-SZ + SPTpol Surveys \\
C & Becker et al. & \texttt{2103.16610} & The mean free path of ionizing photons at 5 < z < 6: evidence for rapid evolution near reionization \\
C & Bosman et al. & \texttt{2108.03699} & Hydrogen reionisation ends by z=5.3: Lyman-$\alpha$ optical depth measured by the XQR-30 sample \\
C & Zhu et al. & \texttt{2109.06295} & Chasing the Tail of Cosmic Reionization with Dark Gap Statistics in the Ly$\alpha$ Forest over 5 < z < 6 \\
C & Villasenor et al. & \texttt{2111.00019} & Inferring the Thermal History of the Intergalactic Medium from the Properties of the Hydrogen and Helium Lyman-alpha Forest \\
C & HERA Collab. et al. & \texttt{2210.04912} & Improved Constraints on the 21 cm EoR Power Spectrum and the X-Ray Heating of the IGM with HERA Phase I Observations \\
C & Fan et al. & \texttt{2212.06907} & Quasars and the Intergalactic Medium at Cosmic Dawn \\
C & Gaikwad et al. & \texttt{2304.02038} & Measuring the photo-ionization rate, neutral fraction and mean free path of HI ionizing photons at 4.9  leq z  leq 6.0 from a larg \\
C & Umeda et al. & \texttt{2306.00487} & JWST Measurements of Neutral Hydrogen Fractions and Ionized Bubble Sizes at z=7-12 Obtained with Ly$\alpha$ Damping Wing Absorptio \\
C & Shao et al. & \texttt{2307.04130} & The 21-cm forest as a simultaneous probe of dark matter and cosmic heating history \\
C & Durovcikova et al. & \texttt{2401.10328} & Chronicling the reionization history at 6 lesssim z  lesssim 7 with emergent quasar damping wings \\
C & Davies et al. & \texttt{2406.18186} & The Predicament of Absorption-dominated Reionization II: Observational Estimate of the Clumping Factor at the End of Reionization \\
C & Sims et al. & \texttt{2504.09725} & Rapid and Late Cosmic Reionization Driven by Massive Galaxies: a Joint Analysis of Constraints from 21-cm, Lyman Line \& CMB Data  \\
C & Ghara et al. & \texttt{2505.00373} & Constraints on the state of the IGM at z sim 8-10 using redshifted 21-cm observations with LOFAR \\
C & Davies et al. & \texttt{2510.25829} & Updated dark pixel fraction constraints on reionization's end from the Lyman-series forests of XQR-30 \\
C & Chaubal et al. & \texttt{2601.20551} & SPT-3G D1: A Measurement of Secondary Cosmic Microwave Background Anisotropy Power \\
C & Soltinsky et al. & \texttt{2607.15341} & Forest without Trees is still Fruitful: Constraints on the thermal state of the neutral IGM at z approx5.6 with the 21-cm forest p \\
C & Montero-Camacho et al. & \texttt{2607.26373} & No way ou$\tau$: Epoch of Reionization Observations Do not Support Large Values of the Optical Depth to Reionization \\
\end{longtable}
